\expandafter\ifx\csname CAPETINCLUDEMODE\endcsname\relax
\documentclass[11pt,a4paper]{article}
% preamble.tex — packages et commandes communes a toutes les fiches CAPET SII IE
\input{glyphtounicode}
\pdfgentounicode=1
\usepackage{cmap}
\usepackage[T1]{fontenc}
\usepackage[utf8]{inputenc}
\usepackage[french]{babel}
% babel[french] active a \begin{document} des raccourcis de ponctuation (:;!?)
% qui cassent la syntaxe to[...] de circuitikz (erreur "+ or - expected") ;
% on les desactive, ce qui ne change que l'espacement fin avant ces signes.
\AtBeginDocument{\shorthandoff{:;!?}}
\usepackage[a4paper, top=1.6cm, bottom=1.6cm, left=1.4cm, right=1.4cm]{geometry}
\setlength{\headheight}{14.5pt}
\usepackage{amsmath, amssymb, amsfonts}
\usepackage{siunitx}
\sisetup{locale=FR, output-decimal-marker={,}, inter-unit-product=\ensuremath{{}\cdot{}},
         group-minimum-digits=4}
% Unites non fournies par siunitx, declarees ici une fois pour toutes (fiches + exercices).
% Les gardes \ifdefined evitent tout conflit avec une redeclaration locale restee dans un
% fichier : ne pas retirer les gardes, ne pas redeclarer ces unites dans les fichiers.
\DeclareSIUnit{\jour}{j}
\ifdefined\var\else\DeclareSIUnit{\var}{var}\fi   % puissance reactive
\ifdefined\vad\else\DeclareSIUnit{\vad}{VAD}\fi   % puissance deformante
\ifdefined\tr\else\DeclareSIUnit{\tr}{tr}\fi      % tour (tr/min)
\usepackage{array, booktabs, multirow, tabularx}
\usepackage{xcolor}
\usepackage[most, breakable]{tcolorbox}
\usepackage{multicol}
\usepackage{enumitem}
\usepackage{mathtools}
\usepackage{lmodern}
% Refonte visuelle 2026 (direction "moderne coloré") : Lato pour le texte,
% maths sans-serif assorties via newtxsf (charge apres amsmath/amssymb).
\usepackage[default]{lato}
\usepackage{newtxsf}
\usepackage{microtype}
\usepackage{fancyhdr}
\usepackage{tikz}
\usetikzlibrary{positioning,arrows.meta}
\usepackage{circuitikz}
\usepackage{lastpage}

% ── Couleurs par domaine (palette modernisée 2026) ────────────────────
\definecolor{colElec}{RGB}{29, 78, 216}        % Électrotechnique — bleu vif
\definecolor{colElecLight}{RGB}{219, 234, 254} % fond bleu clair
\definecolor{colAuto}{RGB}{21, 128, 61}        % Automatique — vert
\definecolor{colMeca}{RGB}{194, 65, 12}        % Mécanique — orange
\definecolor{colTherm}{RGB}{185, 28, 28}       % Thermique — rouge
\definecolor{colQT}{RGB}{232, 244, 233}        % (hérité, plus utilisé par questiontype)
\definecolor{colJury}{RGB}{255, 248, 225}      % fond jury (crème)
\definecolor{colJuryFrame}{RGB}{202, 138, 4}   % bordure/bandeau jury (ambre)
\definecolor{colGrayLight}{RGB}{245, 245, 247} % fond gris neutre
\definecolor{colSeparator}{RGB}{180, 180, 180} % gris séparateur
\definecolor{colExo}{RGB}{8, 145, 178}          % Banque d'exercices — cyan (distinct des fiches)
% ── Encre et gris de la refonte ───────────────────────────────────────
\definecolor{colInk}{RGB}{23, 32, 42}          % texte courant (quasi-noir)
\definecolor{colSlate}{RGB}{83, 98, 118}       % gris ardoise (méta, en-têtes)
\definecolor{colRule}{RGB}{203, 213, 225}      % filets clairs
\definecolor{colJuryBg}{RGB}{255, 248, 225}    % fond crème jury
\definecolor{colPresence}{RGB}{83, 98, 118}    % bloc 5 (présence dans les sujets)
\AtBeginDocument{\color{colInk}}

% ── Style encadré exercice + corrigé (banque d'exercices) ─────────────
\tcbset{
  exoBox/.style={
    breakable,
    enhanced,
    colback=colExo!8!white,
    colframe=colExo,
    colbacktitle=colExo,
    coltitle=white,
    fonttitle=\bfseries\small,
    left=4pt, right=4pt, top=3pt, bottom=3pt,
    before skip=5pt, after skip=5pt,
  }
}

% ── Bascule enonce / corrige (banque d'exercices) ─────────────────────
% \CAPETCORRIGE est defini en ligne de commande (pdflatex "\def\CAPETCORRIGE{1}\input{...}")
% pour compiler la variante corriges d'un meme fichier source ; par defaut (non defini),
% la compilation produit la variante enonces. \CAPETSHOW{enonce}{corrige} choisit la branche.
%
% TROIS MODES (refonte 2026-09-20, manuel de revision) :
%   - aucun drapeau   -> imprime l'enonce            (variante enonces, INCHANGE)
%   - \CAPETCORRIGE   -> imprime le corrige          (variante corriges, INCHANGE)
%   - \CAPETMANUEL    -> imprime l'enonce ET met le corrige de cote pour l'annexe
%
% Le mode manuel stocke chaque corrige dans une macro (\csgdef), sans ecriture
% fichier ni re-tokenisation : le contenu (circuitikz, accents, %) traverse intact.
% L'annexe les rejoue ensuite dans l'ordre via \CAPETANNEXE.
\usepackage{etoolbox}

\newcounter{CAPETexo}          % numero continu d'exercice dans le manuel
\setcounter{CAPETexo}{0}
\newcounter{CAPETannexe}       % compteur de relecture en annexe

% Titre courant de la notion en cours (pose par \CAPETNOTION), sert d'intertitre annexe
\newcommand{\CAPETnotioncode}{}
\newcommand{\CAPETnotiontitre}{}

% \CAPETNOTION{code}{titre} : ouvre une notion dans le manuel (memorise le libelle,
% et marque le rang du premier exercice de la notion pour l'intertitre d'annexe).
\newcommand{\CAPETNOTION}[2]{%
  \renewcommand{\CAPETnotioncode}{#1}%
  \renewcommand{\CAPETnotiontitre}{#2}%
  \csgdef{CAPETnotionstart\the\numexpr\value{CAPETexo}+1\relax}{#1 -- #2}%
}

% ── \degree robuste ───────────────────────────────────────────────────
% Certains fichiers d'exercices utilisent \degree. Ni le preambule ni siunitx
% ne le definissent, et \providecommand local ne suffit pas : sous hyperref,
% \degree existe deja comme commande factice (\ERROR), donc \providecommand
% s'abstient et l'appel echoue a la compilation du manuel (erreur invisible
% en compilation isolee). On le (re)definit donc inconditionnellement ici,
% apres le chargement des packages, pour tout le corpus.
\AtBeginDocument{\providecommand{\degree}{}\renewcommand{\degree}{\ensuremath{^\circ}}}

% Libelle decoratif : les fichiers d'exercices utilisent aussi \CAPETSHOW dans
% leur \fancyhead et leur bloc de titre, pour afficher "(enonces)" ou "(corriges)".
% Ces appels-la ne sont PAS des exercices : en mode manuel ils ne doivent ni
% consommer un numero, ni alimenter l'annexe. On les reconnait a leur premier
% argument, strictement egal au libelle ci-dessous.
\newcommand{\CAPETLIBELLEENONCES}{(\'enonc\'es)}

\ifdefined\CAPETMANUEL
  % --- Mode manuel : enonce en place, corrige differe vers l'annexe ---
  \newcommand{\CAPETSHOW}[2]{%
    \def\CAPETargun{#1}%
    \ifx\CAPETargun\CAPETLIBELLEENONCES
      #1%            appel decoratif : rendu tel quel, hors numerotation
    \else
      \CAPETSHOWEXO{#1}{#2}%
    \fi
  }
  \newcommand{\CAPETSHOWEXO}[2]{%
    \stepcounter{CAPETexo}%
    \csgdef{CAPETcorrige\theCAPETexo}{#2}%
    \edef\CAPETtmpnum{\theCAPETexo}%
    #1%
    \par\vspace{2pt}%
    {\footnotesize\itshape\color{colExo}%
     \hyperref[capet:cor:\CAPETtmpnum]{Corrige de l'exercice \CAPETtmpnum\ $\rightarrow$ p.~\pageref{capet:cor:\CAPETtmpnum}}%
     \label{capet:eno:\CAPETtmpnum}}%
  }
  % \CAPETANNEXE : rejoue tous les corriges captures, dans l'ordre, avec intertitres
  \newcommand{\CAPETANNEXE}{%
    \setcounter{CAPETannexe}{0}%
    \loop\ifnum\value{CAPETannexe}<\value{CAPETexo}%
      \stepcounter{CAPETannexe}%
      \edef\CAPETcurnum{\theCAPETannexe}%
      \ifcsname CAPETnotionstart\CAPETcurnum\endcsname
        \par\medskip
        {\normalsize\bfseries\color{colExo}\csuse{CAPETnotionstart\CAPETcurnum}}%
        \par\nobreak\smallskip
      \fi
      \begin{tcolorbox}[exoBox, breakable,
        title={Corrige de l'exercice \CAPETcurnum\
               \hfill\normalfont\footnotesize
               \hyperref[capet:eno:\CAPETcurnum]{enonce p.~\pageref{capet:eno:\CAPETcurnum}}}]%
      \label{capet:cor:\CAPETcurnum}%
      \csuse{CAPETcorrige\CAPETcurnum}%
      \end{tcolorbox}%
    \repeat
  }
\else
  % --- Modes historiques, comportement strictement inchange ---
  \newcommand{\CAPETSHOW}[2]{\ifdefined\CAPETCORRIGE #2\else #1\fi}
  \newcommand{\CAPETANNEXE}{}
\fi

% ── Difficulté 1 à 3 (étoiles pleines, amssymb) ───────────────────────
\newcommand{\difficulte}[1]{%
  \ifcase#1\relax
  \or $\bigstar$\;niveau 1/3%
  \or $\bigstar\bigstar$\;niveau 2/3%
  \or $\bigstar\bigstar\bigstar$\;niveau 3/3%
  \fi
}

% ── Style encadré extrait de document technique (banque d'exercices) ──
\tcbset{
  docTech/.style={
    breakable,
    enhanced,
    colback=colGrayLight,
    colframe=colSeparator,
    fonttitle=\bfseries\small\itshape,
    coltitle=black,
    colbacktitle=colGrayLight,
    left=4pt, right=4pt, top=3pt, bottom=3pt,
    before skip=4pt, after skip=4pt,
  }
}

% ── Style tcolorbox de section principale ────────────────────────────
% Refonte 2026 : bandeau titre coloré arrondi + fond très pâle, sans cadre.
\tcbset{
  ficheSection/.style={
    breakable, enhanced jigsaw,
    frame hidden, boxrule=0pt, arc=1.8mm,
    colback=#1!4!white,
    colbacktitle=#1,
    coltitle=white,
    fonttitle=\sffamily\bfseries,
    left=7pt, right=7pt, top=5pt, bottom=5pt,
    toptitle=3pt, bottomtitle=3pt,
    before skip=10pt, after skip=8pt,
  }
}

% ── Style encadré question type ──────────────────────────────────────
% Carte blanche, bandeau sombre, liseré très clair.
\tcbset{
  questiontype/.style={
    breakable, enhanced jigsaw,
    colback=white,
    frame hidden, boxrule=0pt, arc=1.8mm,
    colbacktitle=colElec!78!black,
    coltitle=white,
    fonttitle=\sffamily\bfseries\small,
    left=7pt, right=7pt, top=4pt, bottom=4pt,
    toptitle=2.5pt, bottomtitle=2.5pt,
    before skip=8pt, after skip=8pt,
    borderline={0.6pt}{0pt}{colElec!12!white},
  }
}

% ── Style encadré jury ───────────────────────────────────────────────
% Bandeau ambre, fond crème, sans cadre.
\tcbset{
  jury/.style={
    breakable, enhanced jigsaw,
    colback=colJuryBg,
    frame hidden, boxrule=0pt, arc=1.8mm,
    colbacktitle=colJuryFrame,
    coltitle=white,
    fonttitle=\sffamily\bfseries,
    left=7pt, right=7pt, top=4pt, bottom=4pt,
    toptitle=3pt, bottomtitle=3pt,
    before skip=10pt, after skip=8pt,
  }
}

% ── En-tête / pied de page communs (refonte 2026) ─────────────────────
% #1 code fiche, #2 titre court, #3 couleur domaine
\newcommand{\ficheHeaders}[3]{%
  \fancyhead[L]{\small\sffamily{\bfseries\textcolor{#3}{#1}}\hspace{0.7em}\textcolor{colSlate}{#2}}%
  \fancyhead[R]{\small\sffamily\textcolor{colSlate}{CAPET SII IE\, --- \,Session 2026}}%
  \fancyfoot[C]{\small\sffamily\textcolor{colSlate}{\thepage\,/\,\pageref{LastPage}}}%
  \fancyhead[C]{}\fancyfoot[L]{}\fancyfoot[R]{}%
}

% ── Bloc en-tête de fiche (BLOC 0, refonte 2026) ──────────────────────
% #1 code  #2 titre  #3 couleur domaine  #4 domaine  #5 ligne présence  #6 sous-thèmes
\newcommand{\ficheEntete}[6]{%
  \begin{tcolorbox}[enhanced, breakable,
    frame hidden, boxrule=0pt, arc=2.5mm,
    interior style={left color=#3, right color=#3!75!black},
    left=9pt, right=9pt, top=7pt, bottom=7pt,
    before skip=0pt, after skip=6pt]
    {\color{white}\sffamily\bfseries\LARGE #1\quad #2}\\[4pt]
    {\color{#3!22!white}\sffamily\small \MakeUppercase{#4}\quad\textbar\quad #5}\\[1pt]
    {\color{#3!35!white}\sffamily\footnotesize Sous-thèmes : #6}
  \end{tcolorbox}%
}

% ── Mise en page ─────────────────────────────────────────────────────
\pagestyle{fancy}
\fancyhf{}
\renewcommand{\headrulewidth}{0.5pt}
\renewcommand{\footrulewidth}{0pt}
\AtBeginDocument{%
  \renewcommand{\headrule}{{\color{colRule}\hrule width\headwidth height 0.5pt}}%
}

\setlist[itemize]{leftmargin=1.2em, itemsep=1pt, topsep=2pt, parsep=0pt}
\setlist[enumerate]{leftmargin=1.2em, itemsep=3pt, topsep=2pt, parsep=0pt}

\setlength{\parskip}{2pt}
\setlength{\columnsep}{10pt}

\def\CAPETBEGINFILE{\begin{document}}
\def\CAPETENDFILE{\end{document}}
\else
\def\CAPETBEGINFILE{}
\def\CAPETENDFILE{}
\fi

\ficheHeaders{T03}{Complexes et impédances}{colElec}

\CAPETBEGINFILE

% -----------------------------------------------------------------------
%  BLOC EN-TETE
% -----------------------------------------------------------------------

\ficheEntete{T03}{Représentation complexe et impédances}{colElec}%
  {électrotechnique / électronique}%
  {\textbf{Présence :} \textbf{12/13} sujets (2013-2026) - thème \emph{très fréquent}}%
  {impédances R, L, C - loi d'Ohm complexe - déphasage - résonance - modèle MAS - fonctions de transfert}

% -----------------------------------------------------------------------
%  1. RAPPELS DE COURS
% -----------------------------------------------------------------------

\begin{tcolorbox}[
  ficheSection=colElec,
  title={1. Rappels de cours},
]

% -- 1.1 Impedances elementaires --------------------------------------
\textbf{1.1 Impédances des composants élémentaires}

\smallskip
En régime sinusoïdal permanent, on note $\underline{U} = \underline{Z} \cdot \underline{I}$ (loi d'Ohm complexe).
La pulsation est $\omega = 2\pi f$.

\smallskip
\begin{center}
\renewcommand{\arraystretch}{1.45}
\begin{tabular}{llll}
\toprule
\textbf{Composant} & \textbf{Impédance $\underline{Z}$} & \textbf{Module $|\underline{Z}|$} & \textbf{Argument} \\
\midrule
Résistance $R$ & $R$ & $R$ & $0^\circ$ (pas de déphasage) \\[4pt]
Inductance $L$ & $jL\omega$ & $L\omega$ & $+90^\circ$ ($u$ en avance sur $i$) \\[4pt]
Condensateur $C$ & $\dfrac{1}{jC\omega} = \dfrac{-j}{C\omega}$ & $\dfrac{1}{C\omega}$ & $-90^\circ$ ($u$ en retard sur $i$) \\
\bottomrule
\end{tabular}
\end{center}

\medskip
\hrule\vspace{4pt}

% -- 1.2 Associations ------------------------------------------------
\textbf{1.2 Associations RL et RC série - module et déphasage}

\smallskip
\begin{center}
\renewcommand{\arraystretch}{1.45}
\begin{tabular}{lll}
\toprule
\textbf{Circuit} & \textbf{Module} & \textbf{Déphasage $\varphi = \arg(\underline{Z})$} \\
\midrule
RL série & $|\underline{Z}| = \sqrt{R^2 + (L\omega)^2}$ & $\varphi = \arctan\!\left(\dfrac{L\omega}{R}\right) > 0$ (inductif) \\[8pt]
RC série & $|\underline{Z}| = \sqrt{R^2 + (1/C\omega)^2}$ & $\varphi = -\arctan\!\left(\dfrac{1}{RC\omega}\right) < 0$ (capacitif) \\[6pt]
RLC série & $|\underline{Z}| = \sqrt{R^2 + (L\omega - 1/C\omega)^2}$ & $\varphi = \arctan\!\left(\dfrac{L\omega - 1/C\omega}{R}\right)$ \\
\bottomrule
\end{tabular}
\end{center}

\smallskip
\textbf{Diagramme de Fresnel (circuit RL série) :}
\begin{center}
\begin{tikzpicture}[scale=1, font=\small, >=Stealth]
  \draw[->, thin, gray] (-0.3,0) -- (4.4,0);
  \draw[->, thick, colElec] (0,0) -- (3,0) node[below right] {$\underline{I}$};
  \draw[->, thick, colElec!55] (0,0) -- (2.1,0) node[below] {$\underline{U}_R = R\underline{I}$};
  \draw[->, thick, colElec!55] (2.1,0) -- (2.1,1.8) node[midway, right] {$\underline{U}_L = jL\omega\underline{I}$};
  \draw[->, thick, colElec] (0,0) -- (2.1,1.8) node[above right, xshift=2pt] {$\underline{U} = \underline{U}_R+\underline{U}_L$};
  \draw (0.7,0) arc[start angle=0, end angle=40.6, radius=0.7];
  \node at (1.05,0.28) {\footnotesize $\varphi$};
\end{tikzpicture}
\end{center}
$\underline{I}$ est pris comme origine des phases (axe de référence) ; $\underline{U}_R$ est en phase avec $\underline{I}$, $\underline{U}_L$ est en quadrature avance ($+90^\circ$). Pour un circuit RC, $\underline{U}_C$ serait en quadrature retard, plaçant $\underline{U}$ sous l'axe de référence ($\varphi<0$).

\medskip
\hrule\vspace{4pt}

% -- 1.3 Resonance serie ----------------------------------------------
\textbf{1.3 Résonance en série RLC}

À la résonance, $X_L = X_C$ : l'impédance est purement résistive ($\varphi = 0$, courant maximal).

\begin{center}
\renewcommand{\arraystretch}{1.35}
\begin{tabular}{lll}
\toprule
\textbf{Grandeur} & \textbf{Expression} & \textbf{Unité} \\
\midrule
Pulsation propre & $\omega_0 = \dfrac{1}{\sqrt{LC}}$ & rad/s \\[6pt]
Fréquence de résonance & $f_0 = \dfrac{1}{2\pi\sqrt{LC}}$ & Hz \\[6pt]
Facteur de qualité & $Q = \dfrac{L\omega_0}{R} = \dfrac{1}{RC\omega_0}$ & - \\[6pt]
Surtension à résonance & $U_L = U_C = Q \cdot U_{alim}$ & V \\
\bottomrule
\end{tabular}
\end{center}

\medskip
\hrule\vspace{4pt}

% -- 1.4 Fonctions de transfert ---------------------------------------
\textbf{1.4 Formes canoniques des fonctions de transfert}

\smallskip
\begin{center}
\renewcommand{\arraystretch}{1.45}
\begin{tabular}{lll}
\toprule
\textbf{Ordre} & \textbf{Forme canonique} & \textbf{Paramètres} \\
\midrule
1\textsuperscript{er} ordre & $H(j\omega) = \dfrac{H_0}{1 + j\omega/\omega_c}$ & $\omega_c = 1/\tau$ : pulsation de coupure \\[8pt]
2\textsuperscript{e} ordre  & $H(j\omega) = \dfrac{H_0}{1 + 2m\,\dfrac{j\omega}{\omega_0} - \left(\dfrac{\omega}{\omega_0}\right)^2}$ & $\omega_0$ : puls. propre, $m$ : amortissement \\
\bottomrule
\end{tabular}
\end{center}

\smallskip
{\footnotesize $m < 1$ : régime pseudo-périodique. $m \geq 1$ : apériodique. $m \approx 0{,}7$ : optimum ($M_\varphi \approx 65^\circ$).}

\end{tcolorbox}

% -----------------------------------------------------------------------
%  2. FORMULAIRE SYNTHETIQUE
% -----------------------------------------------------------------------

\begin{tcolorbox}[
  ficheSection=colElec,
  title={2. Formulaire synthétique},
]
\small
\begin{multicols}{2}
\renewcommand{\arraystretch}{1.45}
\begin{tabular}{@{}p{3.8cm}ll@{}}
\toprule
\textbf{Formule} & \textbf{Usage} & \textbf{U.} \\
\midrule
\multicolumn{3}{@{}l}{\footnotesize\textit{Impédances élémentaires}} \\
$\underline{Z}_R = R$ & Résistance & $\Omega$ \\
$\underline{Z}_L = jL\omega$ & Inductance & $\Omega$ \\
$\underline{Z}_C = 1/(jC\omega)$ & Condensateur & $\Omega$ \\
\midrule
\multicolumn{3}{@{}l}{\footnotesize\textit{Circuits RL / RC}} \\
$|\underline{Z}_{RL}| = \sqrt{R^2+(L\omega)^2}$ & Module & $\Omega$ \\
$|\underline{Z}_{RC}| = \sqrt{R^2+(1/C\omega)^2}$ & Module & $\Omega$ \\
$\varphi_{RL} = \arctan(L\omega/R)$ & Déphasage & - \\
$\varphi_{RC} = -\arctan(1/RC\omega)$ & Déphasage & - \\
\bottomrule
\end{tabular}

\columnbreak

\renewcommand{\arraystretch}{1.45}
\begin{tabular}{@{}p{3.8cm}ll@{}}
\toprule
\textbf{Formule} & \textbf{Usage} & \textbf{U.} \\
\midrule
\multicolumn{3}{@{}l}{\footnotesize\textit{Résonance série RLC}} \\
$\omega_0 = 1/\sqrt{LC}$ & Puls. propre & rad/s \\
$f_0 = 1/(2\pi\sqrt{LC})$ & Fréq. résonance & Hz \\
$Q = L\omega_0/R$ & Facteur qualité & - \\
\midrule
\multicolumn{3}{@{}l}{\footnotesize\textit{Diagramme de Bode}} \\
$G_{dB} = 20\log|H(j\omega)|$ & Gain & dB \\
$\varphi = \arg H(j\omega)$ & Phase & - \\
$M_\varphi = 180^\circ + \varphi(\omega_{c0})$ & Marge phase & - \\
\midrule
\multicolumn{3}{@{}l}{\footnotesize\textit{Modèle MAS (monophasé équiv.)}} \\
$\underline{I}_r' = \underline{V}_s/(jL_f\omega_s + R_r'/g)$ & Courant rotor & A \\
\bottomrule
\end{tabular}
\end{multicols}
\normalsize
\end{tcolorbox}

% -----------------------------------------------------------------------
%  3. QUESTIONS TYPES
% -----------------------------------------------------------------------

\begin{tcolorbox}[
  ficheSection=colElec,
  title={3. Questions types - extraites des sujets réels},
]

% -- QT1 -------------------------------------------------------------

\begin{tcolorbox}[questiontype,
  title={QT1 - Modèle monophasé équivalent MAS - courant rotorique \hfill\normalfont\footnotesize E2 disciplinaire 2024}
]
\textbf{Énoncé.}
La machine asynchrone est modélisée par son schéma monophasé équivalent.
Exprimer le courant rotorique ramené au stator $\underline{I}_r'$ en fonction de $\underline{V}_s$, $L_f$, $\omega_s$, $R_r'$ et du glissement $g$.

\tcbline

\textbf{Corrigé.}

\textbf{Étape 1 — Impédance du circuit rotorique.}
\[
  \underline{Z}_{rot} = j L_f \omega_s + \frac{R_r'}{g}
\]
(La partie résistive $R_r'/g$ varie avec le glissement : à l'arrêt $g=1$, $\underline{Z}_{rot}=R_r'+jL_f\omega_s$ ; en régime $g\approx0{,}03$, $R_r'/g \approx 33R_r'$.)

\textbf{Étape 2 — Application de la loi d'Ohm complexe.}
\[
  \underline{I}_r' = \frac{\underline{V}_s}{\underline{Z}_{rot}} = \frac{\underline{V}_s}{j L_f \omega_s + R_r'/g}
\]

\textbf{Conclusion.} En régime établi ($g\approx0{,}03$), le terme résistif $R_r'/g$ domine largement le terme inductif : le courant rotorique est quasiment en phase avec $\underline{V}_s$, cohérent avec un fonctionnement proche du synchronisme.
\end{tcolorbox}

\vspace{4pt}

% -- QT2 -------------------------------------------------------------

\begin{tcolorbox}[questiontype,
  title={QT2 - Déphasage d'un filtre - expression analytique \hfill\normalfont\footnotesize Rapport jury 2018 (partie C)}
]
\textbf{Énoncé.}
Un filtre de premier ordre est caractérisé par $H(j\omega)$. Donner l'expression analytique complète de la tension de sortie $v_s(t)$ lorsque l'entrée est $v_e(t) = V_m\cos(\omega t)$.

\tcbline

\textbf{Corrigé.}

\textbf{Étape 1 — Module et argument de la fonction de transfert.}
\[
  |H(j\omega)| = \frac{H_0}{\sqrt{1+(\omega/\omega_c)^2}}
  \qquad
  \varphi(\omega) = \arg H(j\omega) = -\arctan\!\left(\frac{\omega}{\omega_c}\right)
\]

\textbf{Étape 2 — Expression temporelle de sortie.}
\[
  v_s(t) = |H(j\omega)| \cdot V_m \cos\!\bigl(\omega t + \varphi(\omega)\bigr)
\]

\textbf{Conclusion.} Le jury 2018 signale que beaucoup de candidats donnent la valeur numérique de $\varphi$ mais omettent d'écrire l'expression temporelle complète $v_s(t)$ : c'est cette expression complète, pas la seule valeur de $\varphi$, qui est notée.
\end{tcolorbox}

\vspace{4pt}

% -- QT3 -------------------------------------------------------------

\begin{tcolorbox}[questiontype,
  title={QT3 - Marge de phase et stabilité \hfill\normalfont\footnotesize E2 disciplinaire 2024 (partie C)}
]
\textbf{Énoncé.}
La fonction de transfert en boucle ouverte $H_{BO}(j\omega)$ est donnée.
Calculer la marge de phase $M_\varphi$ et conclure sur la stabilité du système.

\tcbline

\textbf{Corrigé.}

\textbf{Étape 1 — Pulsation de coupure à $0\,\text{dB}$.}
Résoudre $|H_{BO}(j\omega_{c0})| = 1$ (ou $G_{dB} = 0$\,dB).

\textbf{Étape 2 — Calcul de la marge de phase.}
\[
  M_\varphi = 180^\circ + \arg\!\bigl(H_{BO}(j\omega_{c0})\bigr)
\]

\textbf{Conclusion.} $M_\varphi > 45^\circ$ : système stable avec dépassement acceptable. $M_\varphi < 0^\circ$ : système instable. \textit{Jury 2024 : seulement 31\,\% des candidats ont traité correctement la marge de phase.}
\end{tcolorbox}

\end{tcolorbox}

% -----------------------------------------------------------------------
%  4. ATTENTES DU JURY
% -----------------------------------------------------------------------

\begin{tcolorbox}[
  jury,
  title={4. Attentes du jury - Erreurs récurrentes et points de vigilance},
]

\begin{enumerate}[label=\textbf{\arabic*.}, leftmargin=*, itemsep=5pt]

  \item \textbf{Oublier $j$ dans l'impédance.}
  Écrire $Z_L = L\omega$ au lieu de $\underline{Z}_L = jL\omega$ supprime toute information de déphasage et fausse les calculs de puissance.

  \item \textbf{Déphasage = argument de $\underline{Z}$, pas de $1/\underline{Z}$.}
  $\varphi = \arg(\underline{Z})$ donne le retard de $i$ sur $u$. Confondre avec l'argument de l'admittance inverse le signe.

  \item \textbf{Marge de phase $\neq$ phase à $0\,\text{dB}$.}
  $M_\varphi = 180^\circ + \arg(H_{BO})$ : ne pas oublier d'ajouter $180^\circ$.
  Le résultat doit être positif pour un système stable.

  \item \textbf{Forme canonique 2\textsuperscript{e} ordre : facteur 2.}
  Le dénominateur est $1 + 2m(j\omega/\omega_0) + (j\omega/\omega_0)^2$.
  Beaucoup de candidats oublient le facteur $2$ devant $m$.

  \item \textbf{Modèle MAS : ne pas confondre $g=1$ (démarrage) et $g\approx0{,}03$ (régime).}
  En régime établi, $R_r'/g \gg jL_f\omega_s$ : le terme inductif est négligeable.
  Au démarrage c'est l'inverse.

  \item \textbf{Expression analytique complète.}
  Une valeur numérique de déphasage n'est pas suffisante : écrire $v_s(t) = \cdots\cos(\omega t + \varphi)$ avec toutes les grandeurs exprimées.

  \item \textbf{Condensateur de démarrage moteur monophasé.}
  Le condensateur $C_p$ crée un déphasage $-90^\circ$ dans l'enroulement auxiliaire, générant un champ tournant elliptique. Ce n'est pas une compensation d'énergie réactive.

\end{enumerate}

\medskip
\begin{tcolorbox}[
  colback=colJuryFrame!12!colJury,
  colframe=colJuryFrame,
  left=4pt, right=4pt, top=2pt, bottom=2pt,
  fontupper=\small\itshape
]
Rapport jury 2024 : «~La partie C (asservissement) a été abordée par seulement 31\,\% des candidats avec correction ; le jury note que la représentation complexe et les diagrammes de Bode font partie des bases attendues à ce niveau.~»\\
Rapport jury 2025 : «~La partie 3 (asservissement) n'a été abordée que par 74\,\% des candidats, et seulement 7\,\% ont obtenu la moyenne sur cette partie.~»
\end{tcolorbox}

\end{tcolorbox}

% -----------------------------------------------------------------------
%  PIED - PRESENCE DANS LES SUJETS
% -----------------------------------------------------------------------

\begin{tcolorbox}[
  ficheSection=colPresence,
  title={Présence dans les sujets analysés (2013-2026)},
]
\begin{center}
\footnotesize\renewcommand{\arraystretch}{1.15}
\begin{tabular}{lccccccccccccc}
\toprule
\textbf{Année} & 13 & 14 & 15 & 16 & 17 & 18 & 19 & 20 & 21 & 22 & 23 & 24 & 25 \\
\midrule
E1 (commune)   & $\checkmark$ & $\checkmark$ & - & $\checkmark$ & $\checkmark$ & $\checkmark$ & $\checkmark$ & $\checkmark$ & $\checkmark$ & $\checkmark$ & $\checkmark$ & $\checkmark$ & $\checkmark$ \\
E2 discipl. IE & - & - & - & - & $\checkmark$ & $\checkmark$ & $\checkmark$ & $\checkmark$ & $\checkmark$ & $\checkmark$ & $\checkmark$ & $\checkmark$ & $\checkmark$ \\
\bottomrule
\end{tabular}
\end{center}
\smallskip
\footnotesize\textbf{Sous-thèmes fréquents :}
impédances R/L/C (11/13) - modèle MAS monophasé (7/13) - fonctions de transfert / Bode (9/13) - résonance LC (5/13) - déphasage analytique (10/13)

\medskip
\textbf{Ordres de grandeur (repères rapides).}
\begin{itemize}[leftmargin=1.3em]
  \item Réseau BT FR : tension phase-neutre typique \SI{230}{\volt}, tension phase-phase typique \SI{400}{\volt}, fréquence \SI{50}{\hertz}.
  \item Courants d'étude fréquents en sujet : de quelques \si{\ampere} à quelques dizaines de \si{\ampere}.
  \item En conversion d'énergie, un rendement attendu reste généralement entre 0{,}8 et 0{,}98 selon la technologie.
\end{itemize}

\medskip
\textbf{Synthèse express (à refaire sans support).}
\begin{enumerate}[leftmargin=1.4em]
  \item Poser les hypothèses du modèle et les conventions de signe.
  \item Écrire 2 à 4 relations clés avec unités.
  \item Réaliser un calcul numérique complet et vérifier l'ordre de grandeur.
  \item Justifier le résultat par un argument technique de niveau CAPET.
\end{enumerate}

\end{tcolorbox}

\CAPETENDFILE
